\section{Modus Ponendo Ponens (MPP)}
\label{69fy}
\lhead{:mathe:logik:kalkül:}

Enthält eine Ableitung eine Formel mit $'\rightarrow'$ als Hauptjunktor als auch
das Antezedenz dieser Formel, dann darf ihr Konsequenz in eine \textit{neue}
Zeile aufgenommen werden. Die Annahmeliste der neuen Zeile wird zusammengesetzt
aus den Annahmelisten der beiden Oberformeln.

\begin{center}
\begin{tabular}{|c|c|c|c|}
  \hline
  Annahme         & Nummer & Formel                & Regel \\
  \hline
  $\alpha$        & (k)    & A                     &  \\
  \hline
  $\beta$         & (l)    & A $\rightarrow$ B     &  \\
  \hline
  $\alpha, \beta$ & (m)    & B                     & l,k MPP \\
  \hline
\end{tabular}
\end{center}

Die Reihenfolge des Vorkommens der Oberformeln ist beliebig.

\textsc{Beispiele:}
\begin{center}
    \[P, P \rightarrow Q \therefore Q\] \\

\begin{tabular}{|c|c|c|c|}
  \hline
  Annahme         & Nummer & Formel             & Regel \\
  \hline
  $\alpha$        & (k)    & $P$                & AE \\
  \hline
  $\beta$         & (l)    & $P \rightarrow Q$  & AE \\
  \hline
  $\alpha, \beta$ & (m)    & $Q$                & l,k MPP \\
  \hline
\end{tabular}
\end{center}

Es gilt: $P, P \rightarrow Q \vdash_J Q$
\newpage
\begin{center}
    \[P \rightarrow (Q\rightarrow R), P \rightarrow Q, P \therefore R\] \\

\begin{tabular}{|c|c|c|c|}
  \hline
  Annahme        & Nummer & Formel                   & Regel \\
  \hline
  1     & (1)    & $P \rightarrow (Q\rightarrow R)$  & AE \\
  \hline
  2     & (2)    & $P \rightarrow Q$                 & AE \\
  \hline
  3     & (3)    & $P$                               & AE \\
  \hline
  1,3   & (4)    & $Q \rightarrow R$                 & 1,3 MPP \\
  \hline
  2,3   & (5)    & $Q$                               & 2,3 MPP \\
  \hline
  1,2,3 & (6)    & $R$                               & 4,5 MPP \\
  \hline
\end{tabular}
\end{center}

Es gilt: $P \rightarrow (Q\rightarrow R), P \rightarrow Q, P \vdash_J R$
\subsection{Referenzen}
\begin{itemize}
    \item \nameref{8hkr} 8hkr.tex
\end{itemize}

\subsection{Literatur}
\begin{itemize}
    \item Timm Lampert - Klassische Logik (2004) Seite 92-93
\end{itemize}
